% from ega1-6-fr.tex
\oldpage[I]{149}
alors, pour toute extension algébriquement close $\Omega$ de $K$, l'espace
sous-jacent à $X\otimes_K\Omega$ a exactement $n$ points, qui s'identifient
aux points de $X$ à valeurs dans $\Omega$.
\end{corollary}

On peut évidemment se borner au cas où l'anneau $A=A(X)$ est local
(\hyperref[I.6.2.2-fr]{6.2.2})~; soient $\mathfrak m$ son idéal maximal,
$L=A/\mathfrak m$ son corps résiduel, extension algébrique de $K$. Les points
de $X$ à valeurs dans $\Omega$ correspondent alors biunivoquement aux
$\Omega$-sections de $X\otimes_K\Omega$
(\hyperref[I.3.4.1-fr]{3.4.1} et \hyperref[I.3.3.14-fr]{3.3.14}), et aussi
aux $K$-homomorphismes de $L$ dans $\Omega$
(\hyperref[I.1.7.3-fr]{1.7.3}), d'où la proposition (Bourbaki, \emph{Alg.},
chap.~V, \textsection7, n\textsuperscript{o}~5, prop.~8), compte tenu de
(\hyperref[I.6.4.3-fr]{6.4.3}).

\begin{env}[6.4.8]
\label{I.6.4.8-fr}
Le nombre $n$ défini dans (\hyperref[I.6.4.7-fr]{6.4.7}) s'appelle le
\emph{rang séparable} de $A$ (ou de $X$) sur $K$, ou aussi le
\emph{nombre géométrique de points de $X$}~; il est donc égal au nombre
d'éléments de $X(\Omega)_K$. Il résulte aussitôt de cette définition que,
pour toute extension $K'$ de $K$, $X\otimes_K K'$ a même nombre géométrique
de points que $X$. Si on désigne ce nombre par $n(X)$, il est clair que si
$X$, $Y$ sont deux schémas finis sur $K$, on a
\[
  n(X\amalg Y)=n(X)+n(Y).
  \tag{6.4.8.1}\label{I.6.4.8.1-fr}
\]

Sous les mêmes hypothèses, on a aussi
\[
  n(X\times_K Y)=n(X)n(Y)
  \tag{6.4.8.2}\label{I.6.4.8.2-fr}
\]
comme il résulte aussitôt de l'interprétation de $n(X)$ comme nombre
d'éléments de $X(\Omega)_K$ et de la formule
(\hyperref[I.3.4.3.1-fr]{3.4.3.1}).
\end{env}

\begin{proposition}[6.4.9]
\label{I.6.4.9-fr}
Soient $K$ un corps, $X$, $Y$ deux $K$-préschémas algébriques, $f:X\to Y$ un
$K$-morphisme, $\Omega$ une extension algébriquement close de $K$, de degré
de transcendance infini sur $K$. Pour que $f$ soit surjectif, il faut et il
suffit que l'application $X(\Omega)_K\to Y(\Omega)_K$ correspondant à $f$
(\hyperref[I.3.4.1-fr]{3.4.1}) soit surjective.
\end{proposition}

La nécessité résulte de (\hyperref[I.6.3.10-fr]{6.3.10}), en remarquant que
$f$ est nécessairement de type fini (\hyperref[I.6.3.9-fr]{6.3.9}). Pour voir
que la condition est suffisante, on raisonne comme dans
(\hyperref[I.6.3.10-fr]{6.3.10}), en remarquant que pour tout $y\in Y$,
$k(y)$ est une extension de $K$ de type fini, et par suite est $K$-isomorphe
à un sous-corps de $\Omega$.

\begin{remark}[6.4.10]
\label{I.6.4.10-fr}
Nous verrons au chap.~IV que la conclusion de
(\hyperref[I.6.4.9-fr]{6.4.9}) est encore valable sans hypothèse relative au
degré de transcendance de $\Omega$ sur $K$.
\end{remark}

\begin{proposition}[6.4.11]
\label{I.6.4.11-fr}
Si $f:X\to Y$ est un morphisme de type fini, pour tout $y\in Y$, la fibre
$f^{-1}(y)$ est un préschéma algébrique sur le corps résiduel $k(y)$ et pour
tout $x\in f^{-1}(y)$, $k(x)$ est une extension de type fini de $k(y)$.
\end{proposition}

Comme $f^{-1}(y)=X\otimes_Y k(y)$
(\hyperref[I.3.6.3-fr]{3.6.3}), la proposition résulte de
(\hyperref[I.6.3.4-fr]{6.3.4, (iv)}) et de
(\hyperref[I.6.3.3-fr]{6.3.3}).

\begin{proposition}[6.4.12]
\label{I.6.4.12-fr}
Soient $f:X\to Y$, $g:Y'\to Y$ deux morphismes~; posons
$X'=X\times_Y Y'$ et soit $f'=f_{(Y')}:X'\to Y'$. Soit $y'\in Y'$,
$y=g(y')$~; si la fibre $f^{-1}(y)$ est un schéma algébrique fini sur $k(y)$,
alors la fibre $f'^{-1}(y')$ est un schéma algébrique fini sur $k(y')$, ayant
même rang et même nombre géométrique de points que $f^{-1}(y)$.
\end{proposition}

Compte tenu de la transitivité des fibres
(\hyperref[I.3.6.5-fr]{3.6.5}), cela résulte aussitôt de
(\hyperref[I.6.4.6-fr]{6.4.6}) et
(\hyperref[I.6.4.8-fr]{6.4.8}).

\begin{env}[6.4.13]
\label{I.6.4.13-fr}
La prop. (\hyperref[I.6.4.11-fr]{6.4.11}) montre que les morphismes de type
fini correspondent intuitivement aux «~familles algébriques de variétés
algébriques~», les points de $Y$ jouant
